\PassOptionsToPackage{table,xcdraw}{xcolor}
\documentclass[11pt, a4paper]{article}

\usepackage[T1]{fontenc}
\usepackage[utf8]{inputenc}
\usepackage{tgpagella}
\usepackage[spanish, es-noshorthands]{babel}
\usepackage{amsmath, amssymb}
\usepackage[most]{tcolorbox}
\usepackage[margin=2cm]{geometry}
\usepackage{fancyhdr}

\pagestyle{fancy}
\fancyhf{}
\setlength{\headheight}{16pt}
\lhead{\textbf{UCB Sede Tarija} | Ing. Mecatrónica}
\rhead{IMT-221 | Práctica Individual W09S01}
\renewcommand{\headrulewidth}{0.4pt}
\definecolor{ucbblue}{RGB}{0, 51, 102}

\begin{document}

\begin{tcolorbox}[colback=ucbblue!5, colframe=ucbblue, title=\textbf{Práctica Individual Tipo Examen: Consolidación Hito 2}]
    Resuelva individualmente detallando sus pasos y suposiciones. \textbf{No use derivaciones largas memorizadas}, aplique las leyes de circuitos directamente[cite: 7].
\end{tcolorbox}

\textbf{Nombre Estudiante:} \hrulefill \quad \textbf{Fecha:} \hrulefill

\vspace{0.3cm}
\section*{Round 1: Preparación (DC y Ganancia Diferencial)[cite: 7]}

\noindent \textbf{Problema A — Q-point[cite: 7]} \\
Dado $V_{CC}=9.0\,\text{V}$, $R_C=4.7\,\text{k}\Omega$, $v_C=6.55\,\text{V}$ y $v_E=-0.70\,\text{V}$[cite: 7]. 
Calcule explícitamente $I_C$ y $V_{CE}$[cite: 7]. Determine la región de operación y evalúe la compatibilidad con el bias nominal de Hito 2[cite: 7].
\vspace{3cm}

\noindent \textbf{Problema B — Current Mirror[cite: 7]} \\
Dado $R_{REF}=8.2\,\text{k}\Omega$, $V_{EE}=-9\,\text{V}$ y $V_{BE}\approx0.7\,\text{V}$[cite: 7]. 
Estime analíticamente $I_{REF}$[cite: 7]. Si en hardware se mide una corriente de cola $I_T=0.88\,\text{mA}$, cuantifique la discrepancia, liste causas físicas y juzgue si esto, por sí solo, demuestra una falla catastrófica del circuito[cite: 7].
\vspace{3cm}

\noindent \textbf{Problema C — Differential Pair ($g_m$ y $A_d$)[cite: 7]} \\
Considerando un punto de reposo validado con $I_C=0.48\,\text{mA}$, y $R_C=4.7\,\text{k}\Omega$[cite: 7].
Calcule $g_m$ (use $V_T=25.85\,\text{mV}$)[cite: 7]. Calcule la ganancia Single-Ended $A_{d,se}$[cite: 7]. Para un $v_d = 5\,\text{mV}$, calcule la amplitud esperada $v_{C2,ac}$ e indique su signo o fase relativa respecto a $v_d$[cite: 7].
\vspace{3.5cm}

\newpage
\section*{Round 2: Cierre (Modo Común, CMRR y Diagnóstico)[cite: 7]}

\noindent \textbf{Problema D — Descomposición Diferencial/Common Mode[cite: 7]} \\
Usted inyecta voltajes asimétricos: $v_1=1.020\,\text{V}$ y $v_2=0.980\,\text{V}$[cite: 7]. 
Calcule analíticamente $v_d$ y $v_{cm}$, clasifique la naturaleza de la excitación y demuestre matemáticamente la reconstrucción de $v_1$ y $v_2$[cite: 7].
\vspace{3cm}

\noindent \textbf{Problema E — CMRR Validity Gate[cite: 7]} \\
Un grupo reporta el \textbf{Caso 1}: $A_{d,se}=40$, $A_{c,se}=0.04$, ambos medidos en $v_{C2}$[cite: 7]. 
Otro grupo reporta el \textbf{Caso 2}: $A_{d,diff}=80$ (entre colectores), $A_{c,se}=0.04$ (en un colector)[cite: 7]. 
Determine rigurosamente la validez del CMRR para cada caso y calcúlelo en dB solo si el método es válido[cite: 7].
\vspace{3cm}

\noindent \textbf{Problema F — Measurement Interpretation[cite: 7]} \\
Un grupo registra la siguiente tabla de datos[cite: 7]:
\begin{itemize}
    \item $v_d$ inyectado = $5.0\,\text{mV}_{pp}$
    \item Salida $v_{C2}$ diferencial = $200\,\text{mV}_{pp}$
    \item $v_{cm}$ inyectado = $200\,\text{mV}_{rms}$
    \item Salida $v_{C2}$ común = $5\,\text{mV}_{pp}$
\end{itemize}
¿Se puede calcular el CMRR directamente operando estos números? Justifique su respuesta detallando el error metodológico[cite: 7].
\vspace{2.5cm}

\end{document}
